\scanpage{91}
5) \uline{Étude par cas}.

I) $\Delta(\alpha)\ne0$ \quad i.e. $\alpha\ne\pm2$ \quad (``cas ordinaire'')

Étendant au besoin le corps de base $k$ en $k(\unsure{\zeta|\rho})=k(\zeta')$, et rapportant maintenant $E$ au nouveau syst. de coordonnées de centre $o$, ayant comme axes les droites propres de $u$ $D_\zeta$ et $D_{\zeta'}$, et pour lesquels $s_0$ est de coordonnées $(1,1)$, on trouve par un argument déjà fait le cas-type $\Pi_\zeta$ (où $\zeta$ n'est pas \unsure{néc$^{t}$} une racine de 1).
Il y a une fonction quadratique (\unsure{p.r. au centre|par son centre} $o$) invariante, canonique à un scalaire multiplicatif près, qu'on peut normaliser en exigeant que sa valeur en les sommets de $\Pi_\zeta$ soit égale à 1.. L'image \deleted{$\Pi_\zeta$} \inserted{de la représentation $\Pi_\zeta$} est un polygône géométrique, dont le nb de \unsure{côté|cté|côtés}, égal à l'ordre de $\zeta$, est (\deleted{s'il} \inserted{quand il} est fini) premier à la caractéristique.

\deleted{II) \uline{$\alpha=-2$}} \deleted{On supposera ici $\alpha\ne2$ (i.e. $p\ne2$),}
\[
\begin{aligned}
&\deleted{\unsure{\sigma|v|r}_0(x,y)=(1-x-2y,y)}\\
&\deleted{u_0(x,y)=(1-y-2x,x)}\\
&\deleted{u_0^2(x,y)=(1-x-2(1-y-2x),1-y-2x)}\\
&\qquad\deleted{=-1+3}
\end{aligned}
\]
\deleted{de sorte que l'on a encore \unsure{l'existence|existence} et le centre $o$.}
% archive sheet number: 90
